% Pista ana-25i-q1 · Anàlisi
% Procedència: PAU Catalunya 2025 · Convocatòria de juny (incidències) · Sèrie 4 · Exercici 1
\documentclass[11pt,a4paper]{article}
\usepackage{pista}
\pistainfo{Catalunya 2025}{Convocatòria de juny (incidències)}{Sèrie 4}

\begin{document}

\section*{\textcolor{blauPista}{Exercici 1}}

\noindent Considereu la funció
\[
f(x) = +\sqrt{1 + x^{3}}.
\]

\subsection*{\textit{a)} \normalfont Indiqueu el domini d'aquesta funció i els talls de la corba $y = f(x)$ amb els eixos de coordenades. \hspace*{\fill}\textnormal{[0,5~punts]}}

\pista{Segueix l'ordre de l'estudi d'una funció. Aquí només et demanen dos passos:
\begin{itemize}
\item \textit{Domini.} L'expressió de dins l'arrel quadrada ha de ser no negativa: estudia quins valors de $x$ compleixen
\[
1 + x^{3} \geq 0.
\]
\item \textit{Talls amb els eixos.} Amb l'eix d'abscisses, resol $f(x) = 0$. Amb l'eix d'ordenades, només cal avaluar $f(0)$.
\end{itemize}}

\subsection*{\textit{b)} \normalfont Resoleu l'equació $f'(x) = 0$ i estudieu les zones de creixement i decreixement de la funció $f(x)$, així com els seus màxims i mínims locals. \hspace*{\fill}\textnormal{[1~punt]}}

\pista{Per derivar, tingues present que $f(x)$ és una composició de funcions.
\begin{itemize}
\item Resol $f'(x) = 0$.
\item Estudia el signe de $f'(x)$ als intervals resultants per classificar els punts trobats.
\end{itemize}

\medskip
\textbf{Atenció:} no t'oblidis dels valors del domini on $f'(x)$ podria no estar definida, especialment a la frontera del domini. I un punt amb $f'(x) = 0$ no sempre és un màxim o un mínim: de vegades és un punt d'inflexió.}

\subsection*{\textit{c)} \normalfont Trobeu quin ha de ser el valor de $a$ per tal que el punt $(a, 3)$ pertanyi a la gràfica de la funció i calculeu l'equació de la recta tangent a la corba $y = f(x)$ en aquest punt. \hspace*{\fill}\textnormal{[1~punt]}}

\pista{Que un punt pertanyi a la gràfica vol dir que les seves coordenades compleixen \mbox{$y = f(x)$}: imposa aquesta condició per trobar $a$.

\medskip
Després, aplica la fórmula de la recta tangent:
\[
y - f(a) = f'(a)\,(x - a)
\]}

\end{document}
